\documentclass{article}
\usepackage{graphicx}
\usepackage{booktabs}
\usepackage{amsmath}
\usepackage{hyperref}

\title{AnnotateROI: An Enterprise Framework for Measuring Annotation Return on Investment and Determining Labeling Sufficiency}
\author{[Intern Name], Spurthi Setty \\
Anote AI}
\date{}

\begin{document}

\maketitle

\begin{abstract}
\noindent
Enterprises investing in data annotation face two recurring questions: how much annotation is enough, and when does the cost of additional labels exceed the marginal performance gain? We present \textbf{AnnotateROI}, a framework that combines (1) a parameterized annotation cost model, (2) a data-driven stopping criterion based on diminishing returns detection, and (3) an LLM-vs-human decision tree that recommends annotation strategy given task type and budget. Building on experimental results from AnnotateBench \cite{annotatebench}, we validate AnnotateROI across five NLP task types and show that our stopping criterion halts annotation at the cost-optimal point with X\% accuracy, reducing over-annotation by an average of Y\%.
\end{abstract}

\section{Introduction}
% The annotation ROI problem: enterprises often over- or under-annotate
% No principled framework for knowing when to stop
% This paper: cost model + stopping criterion + decision tree

\section{Related Work}
\subsection{Active Learning Stopping Criteria}
\subsection{Cost-Sensitive Learning}
\subsection{Data Valuation}

\section{The AnnotateROI Framework}
\subsection{Annotation Cost Model}
% Parameters: hourly rate, examples/hour, LLM API cost, overhead
\begin{equation}
\text{ROI}(n) = \frac{\Delta \text{Accuracy}(n)}{\text{Cost}(n)}
\end{equation}

\subsection{Stopping Criterion}
% Marginal gain detection: stop when ROI(n) < threshold
% Algorithm: sliding window over learning curve slope

\subsection{LLM-vs-Human Decision Tree}
% Decision factors: task complexity, domain expertise required, budget, timeline

\section{Experimental Setup}
% Builds on AnnotateBench (T2a) results
% 3 enterprise budget scenarios: small ($500), medium ($5K), large ($50K)

\section{Results}
% Table: Stopping criterion accuracy by task type
% Figure: ROI curves across annotation strategies
% Figure: Decision tree validation

\subsection{Stopping Criterion Validation}
\subsection{Enterprise Budget Simulation}
\subsection{Sensitivity Analysis}

\section{Conclusion}

\section*{References}
\small
\begin{enumerate}
    \item [T2a] AnnotateBench: A Systematic Benchmark of Annotation Strategies Across NLP Task Types.
    \item Settles, Burr. Active Learning Literature Survey, 2009.
    % Add more references
\end{enumerate}

\end{document}
